(a) The leading digit $a$ is a nonzero remainder modulo $2$, so $a=1$ and the last digit $c$ is odd. Since $c$ is also a remainder modulo $4$, $c=1$ or $3$. The remainder modulo $5$ is then $d=c$. But the last two digits are $11c$, and the condition modulo $4$ requires $11c\equiv c$, or $2c\equiv0\pmod4$. This forces $c$ even, a contradiction. No original integer exists.

(b) Again $a=1$, and $e$ is odd with $0\le e<6$, so $e\in\{1,3,5\}$. The other remainders are forced:
\[
d\equiv e\pmod5,\quad
c\equiv10d+e\pmod4,\quad
b\equiv e\pmod3.
\]
Using the least nonnegative remainders gives the candidates
\[
11311,\quad10133,\quad12105.
\]
Their digit sums are $7,8,9$. Only the first has its required remainder $b$ modulo $3$; it also has all the required other remainders. Thus the corrected integer is $\boxed{11311}$ and its last-two-digit remainder is $\boxed{11}$.
